\subsection{An executed finite composition table with nonzero Wall class}
\label{D:final-table}

The following table is specified by integer indices and total lookup
formulas. Its size makes a rectangular printed array impractical; the
index formulas specify every entry, including every typed composite.
There is no unspecified triangulation or search for a subdivision in
this definition.

\begin{theorem}[An explicit negative answer to D1]
\label{D:Q8-final-table}
There is a connected finite category $\mathcal W_8$ with
\[
\begin{split}
 |\operatorname{Ob}\mathcal W_8|
   &=489\,754\,140\,643\,145\,442,\\
 |\operatorname{Mor}\mathcal W_8|
   &=337\,126\,384\,361\,511\,151\,079\,802.
\end{split}
\]
Its classifying space is finitely dominated and has a CW model of
dimension at most eight. Its fundamental group is $Q_8$, and
\[
 w(B\mathcal W_8)=[\operatorname{im}e_8]\ne0
 \quad\text{in }\widetilde K_0(\mathbb ZQ_8).
\]
Consequently $B\mathcal W_8$ is not homotopy equivalent to a finite
CW complex. The complete indexed table is implemented in
\path{code/d_final_wall_category.py}.
The function \texttt{build\_final\_category()} returns it as
\texttt{Final}; its four operations are
\texttt{source}, \texttt{target}, \texttt{identity}, and
\texttt{compose}.
\end{theorem}

\begin{proof}
We give the finite inputs, all table constructors, and the homotopy and
chain calculations. The implementation is a literal indexing of these
formulas, rather than a criterion used to infer a homotopy type from
a finite homology range.

\paragraph{The group, base, and projector.}
Number $Q_8$ as $2a+s$, where $a=0,1,2,3$ denotes $1,i,j,k$ and
$s=0,1$ is the sign. Multiply signs modulo two; equal nonzero axes
have product $-1$, and $(i,j),(j,k),(k,i)$ have respectively products
$k,i,j$, with the reversed products negative. This specifies all
64 products. Let $L$ be the boundary of the four-dimensional cross
polytope with these eight vertices, the opposite pairs being
$\{g,-g\}$. Left multiplication acts freely on $L$.
Choose the lexicographically least face in each orbit. The category
$K_8$ has these ten representatives as objects and arrows
\[
 (F,H,g):F\longrightarrow H\qquad(gF\subseteq H).
\]
Its identity is $(F,F,1)$ and its composition is
\[
 (H,J,h)(F,H,g)=(F,J,hg).
\]
There are 68 arrows. Its nerve is the orbit of the face-poset nerve
of $L$; hence $BK_8\cong|\operatorname{sd}L|/Q_8$ and has
fundamental group $Q_8$ and universal cover $S^3$. Its nondegenerate
simplex counts are $(10,58,96,48)$ in dimensions zero through three.

Put $R=\mathbb ZQ_8$, $N=\sum_g g$, and use the unimodular change
of basis
\[
 P=\begin{pmatrix}2&5\\1&2\end{pmatrix},\qquad
 e=P^{-1}e_8P=
 \begin{pmatrix}-4&N-10\\N+2&5\end{pmatrix},\qquad q=1-e.
\]
Direct multiplication gives
$eq=qe=(8N-N^2)I=0$ and $e^2=e$. Thus
$[\operatorname{im}e]=[\operatorname{im}e_8]\ne0$ by
Proposition~\ref{D:Q8-projector}.
Each of $e,q$ is a sum of 35 signed matrix units with group labels.

\paragraph{Finite coefficient sets and their exact maps.}
Let $J_n$ have vertices $v_0,\ldots,v_n$ and edges
$u_1,\ldots,u_n$, with the two nonidentity arrows
$v_{i-1}\to u_i\leftarrow v_i$. Put
$A_n=NJ_n/\{v_0,v_n\}$. The quotient identifies the two endpoint
constant simplicial subsets with the basepoint. In degree $k$ it has
$2n(k+1)$ simplices, including the basepoint; monotone pullback acts
on its vertex strings. Take $S=A_1^{\wedge5}$ and
$Y=|NK_8|\vee S\vee S$, using the actual finite simplicial nerve
of $K_8$ as the base and literal pointed simplicial quotients for
the spheres. All spaces below are finite simplicial sets with this
base as a simplicial subset.

A group label $g$ uses the vertex path $(1,g)$, except that $1$ uses
$(1,1)$ and $-1$ uses $(1,i,-1)$. Its interval face category maps
to $K_8$ by face-orbit normalization. A label map from sphere $j$
to sphere $i$ is represented by the cospan whose $i$-th target sphere
has its attachment point transported along this interval. The left
leg uses its initial endpoint, the right leg its final endpoint, and
all unused source spheres collapse to the base. The right leg is a
weak equivalence since the transport interval is contractible.
A negative sign reflects suspension coordinate four, numbered from
zero. All additive pinches initially use coordinate zero, so signed
elementary maps commute strictly with these pinches.

The full endpoint functors, including their projections to $K_8$,
are given by the following finite rule. For a path or disk face $F$
with labelled vertices let $\bar F$ be its least orbit representative,
and choose the least $t_F\in Q_8$ satisfying $t_F\bar F=F$.
Send an inclusion $F\subseteq H$ to
$(\bar F,\bar H,t_H^{-1}t_F)$. Identity and composition follow by
cancellation of the adjacent $t$'s. Each multiplication disk has
boundary the concatenated paths for $h$, $g$, and $(hg)^{-1}$,
and is coned to the least cross-polytope vertex whose opposite is
absent from that boundary. The resulting face-poset disk and its
projection use the same rule. There are exactly 64 fixed disks.
The inner label is used first: an inner $h$ followed by an outer
$g$ uses the disk $(h,g)$.

We spell out the finite homotopy operations combining these data.
The map $A_n\to A_1$ retaining only its last edge has degree one,
while collapsing all subdivision vertices gives
$A_n\to\bigvee_n A_1$. Together they give the nary pinch roof.
For retained-edge collapses $q_a,q_b:J_n\to J_1$, assume $a<b$.
There is a common upper functor $h_{ab}$: vertices below $a$ go
to the left endpoint, vertices from $a$ through $b-1$ to the
edge, and vertices at least $b$ to the right endpoint; edges from
$a$ through $b$ go to the edge, and the remaining edges to the
corresponding endpoint. Thus
$q_a\Rightarrow h_{ab}\Leftarrow q_b$, fixing both endpoints.
Apply these functors independently to branches of an
$n$-by-$n$ suspension grid. Moving all retained edges at once,
reinterpreting the resulting sparse matching by columns, and
reversing the procedure gives any prescribed permutation in four
prism stages. The same construction compares addition in
coordinates zero and four in two prism stages.

Successive pinches of arities $n,m$ are compared to the flat
$nm$ pinch by retaining edges $m,2m,\ldots,nm$ and splitting
$J_{nm}$ into its $n$ consecutive blocks of $m$ edges. These maps
satisfy literally
\[
 q_ns=q_{nm},\qquad p_ns=(\bigvee q_m)t,\qquad
 (\bigvee p_m)t=p_{nm}.
\]
A three-edge mapping cylinder realizes the comparison without
inverting a simplicial map. Finally adjacent opposite terms fold
their two edges into one based $J_1$; the pointed contraction
$\operatorname{id}\Rightarrow\mathrm{peak}\Leftarrow\mathrm{constant}$
contracts every pair simultaneously. A zero tail is sent directly
to the common basepoint. Different coefficient apices are kept in
a target-star mapping cylinder, so this operation never assumes a
strict inverse to a weak equivalence.

All mapping cylinders have explicit simplices: besides endpoint
simplices, a nonconstant cylinder simplex is a pair $(x,c)$ with
$1\leq c\leq k$, representing the monotone map $[k]\to[1]$ whose
first $c$ values are zero. Pullback counts how many selected
vertices lie before $c$; zero or $k+1$ selected zeros use the
specified endpoint map. Relative cylinders identify the base
cylinder with the base. This gives an exact finite normal form
and its pullback function. A chain of cylinders uses an additional
finite factor index. Products, pointed quotients, and wedges act
coordinatewise, retaining the summand tag unless the simplex
hits the base. These are the complete simplicial lookup rules
in the coefficient compiler.

Expanding the two 35-term matrices gives 1225 elementary products
for each order. A mismatched matrix index gives zero; a matched
index uses its fixed group disk. Sorting by matrix entry and group
label, with positive and negative terms adjacent, gives 294 inverse
pairs and 637 zero terms for $eq$, and 300 inverse pairs and 625
zeros for $qe$. The preceding fixed operations give exactly 50
coefficient roofs and 49 actual weak comparison maps for each
nullhomotopy. Every comparison preserves both true $Y$ endpoint
maps. All of this finite input has been executed, including the
two orders and all group/sign/index combinations.

Take the relative mapping-cylinder homotopy colimit $Z_{eq}$ or
$Z_{qe}$ of the corresponding 50-term zigzag. It contains its
initial composite apex and is weakly equivalent to $Y$. The left
legs give an actual map
\[
 H:(Y\times I)/((k,t)\sim k)\longrightarrow Z,\qquad
 H|_0=f'c',\quad H|_1=i_Zr_Y.
\]
Here $I$ is the finite oriented path supplied by the zigzag.
The base collapse is justified: before collapse $K\times I$ is
a simplicial subcomplex of the ordinary mapping cylinder.
Its projection to $K$ is a homotopy equivalence, and pushout along
that subcomplex inclusion preserves the equivalence.
Both endpoint inclusions of the relative domain are weak
equivalences. All coefficient apices have dimension at most
seven; this last path cylinder has dimension at most eight.

\paragraph{The finite category table constructors.}
For each coefficient simplicial set $X$ use the opposite truncated
simplex category
\[
 O_8(X)=(\Delta_{\leq8}\mathbin{\downarrow}X)^{\rm op}.
\]
An object is $(k,x)$, $0\leq k\leq8$, $x\in X_k$. An arrow
from $(k,x)$ to $(l,y)$ is a monotone map
$\theta:[l]\to[k]$ with $y=\theta^*x$. Identity is the identity
ordinal map; composition is composition of ordinal maps. Its
object and arrow counts are respectively
\[
 \sum_{k=0}^8|X_k|,\qquad
 \sum_{k,l=0}^8|X_k|\binom{k+l+1}{l+1}.
\]
Order objects by degree and then by the simplex normal forms;
order outgoing arrows by target degree and the lexicographic
ordinal map. These rules give unique integer indices. The
truncated-simplex comparison of Lemma~\ref{lem:B-prelim} gives
$BO_8(X)\simeq|X|$. To recall the relevant proof, in the comma
category from a simplex to the truncated simplex category,
ordered-image factorization retracts to the surjective ordinal
maps; the unique nondegenerate ancestor is terminal there and
has dimension at most eight. Quillen's Theorem~A compares the
truncated and full simplex categories. The full simplicial
replacement has the homotopy type of $|X|$. Passing to the
opposite category preserves realization. Our convention makes
the number of arrows out of an object depend only on its degree.

Products use ordered pairs and coordinatewise composition.
For a span $A\xleftarrow{F}D\xrightarrow{G}B$, the Grothendieck
category has the three tagged object sets $D,A,B$. Besides
internal arrows it has crossing arrows $(x,a)$ from $x\in D$
to $a_1\in A$ for $a:F(x)\to a_1$, and analogous $B$ arrows.
Composition after a crossing is target composition; composition
before a crossing replaces $(x,a)$ and $u:x'\to x$ by
$(x',aF(u))$. There are no two successive crossings.
The same formulas apply to the four crossings of the two
five-stage pinch diagrams. Identity and associativity therefore
follow from the input laws and functor equations. This is a
universal proof of the table laws, independent of numerical checks.

For exact indexing, partition objects into the finite types used
by these constructors. If a type contains $n_t$ objects and has
$a_t$ outgoing arrows per object, enumerate its objects
consecutively and each object's outgoing arrows in the displayed
internal/crossing order. The block has $n_ta_t$ arrows.
Product types are pairs; span types add their stage tag.
On a crossing, the target arrow is enumerated in the outgoing
block of the image of the source under the given functor.
Decoding these finite mixed-radix indices and applying the
composition rule gives a total lookup for every typed pair.
The table for $\mathcal W_8$ has 3897 such types.

\paragraph{The final finite construction.}
Apply the relative tensor of Lemma~\ref{D:finite-coefficient-tensor}
to the coefficient objects with the two nine-element monoids
$M_0,M_1$. Its two finite Grothendieck spans are strictly
functorial. Form the relative wedge $V$ of the two $Y$ tensors.
Let $W$ be the wedge of the $q$-apex tensor in parity zero and
the $e$-apex tensor in parity one, and let $Z$ be the wedge of
the $Z_{eq},Z_{qe}$ tensors. The initial three factors of each
composite coefficient chain give actual functors
\[
 c':V\longrightarrow W,\qquad f':W\longrightarrow Z.
\]
There are weak right-leg functors $V\to W$ and $V\to Z$.
Under them $c'$ has coefficients $q,e$ and $f'$ has coefficients
$e,q$. The latter assertion uses the explicit junction cylinder
between the two coefficient blocks; its two endpoint maps are
not asserted literally equal.

Tensor the actual $H$ above and form the relative cones
$C=\operatorname{Cof}_K(f')$, $C_{fc}$, and $C_H$.
The endpoint inclusions give equivalences
\[
 C_{fc}\longrightarrow C_H\longleftarrow C_{\mathrm{tensor}(0)}.
\]
The tensor of a coefficient zero is not literally the section
followed by retraction. It factors through $T_K(M_\epsilon)$.
In that base tensor a finite middle functor moves the outer
product stage into the central boundary stage, and has natural
transformations to both the identity and the section--retraction
functor. On the boundary span the middle functor fixes the
product stage and sends its two copies of $K$ to the central
copy. The transformation components are its canonical span
arrows. Both transformations are over $K$ and identity on the
chosen section. Composing with the coefficient retraction and
section gives the further cone equivalences
\[
 C_{\mathrm{tensor}(0)}\longrightarrow C_{\mathrm{middle}}
 \longleftarrow C_0.
\]
For $a\Rightarrow b$, the cone comparison runs
$\operatorname{Cone}(b)\to\operatorname{Cone}(a)$; on a crossing
it sends $t:b(x)\to y$ to $t\eta_x:a(x)\to y$.
These are the directions used here.

The suspension-summand inclusion $\Sigma_KV\to C_0$, these
comparisons, and the strict square $(c',1)$ from $C_{fc}$ to $C$
give a finite cospan representing $\ell:\Sigma_KV\to C$.
A second represents $\psi=jf^c\sigma$. Here $\sigma$ uses the
31-arrow parity shifts of Theorem~\ref{D:parity-shifts}, the
natural finite suspension--tensor comparison, and the relative
wedge; $f^c$ is the complementary coefficient roof.
The implementation retains all base paths. Composition uses
a span homotopy colimit, appending its two canonical natural
transformations to the base path.

The intrinsic pinch uses the height-one diagrams
\[
 K\longleftarrow V\longrightarrow V\longleftarrow V
       \longrightarrow K,\qquad
 K\longleftarrow V\longrightarrow K\longleftarrow V
       \longrightarrow K.
\]
In the first the two middle arrows are identities; all remaining
arrows are retractions. Collapse its first half to obtain a weak
map to $\Sigma_KV$, and retract its middle $V$ to obtain a map
to the second diagram. The ordinary wedge of two suspensions
maps weakly to the second diagram, reversing the ends of the
second suspension. This gives a finite pinch cospan with branch
degrees $(1,-1)$. The wedge/fold construction, with collars on
the stored base paths, gives
\[
 \Sigma_KV\xrightarrow{g_l}G\xleftarrow[\simeq]{g_r}C
\]
representing $g=\ell-\psi$ in this convention.
Every map is strictly over $K$; base paths are retained when a
section is preserved only up to homotopy.
Define $\mathcal W_8$ as the finite Grothendieck construction
of $G\xleftarrow{g_l}\Sigma_KV\to K$. Since $g_r$ is a weak
equivalence over $K$, its realization is the required relative
homotopy cofiber of $g$.

\paragraph{Universal-cover chains and the obstruction.}
The relative universal-cover complex of $V$ is chain-homotopy
equivalent to one free $R^2$ in every degree $n\geq7$, with
zero differential. The coefficient $S^5$ supplies the shift
by five; the parity generators start in degrees two and three.
On this model $f_n$ alternates $e,q$, while its complement
alternates $q,e$. The coefficient nullhomotopies have zero
degree-six cellular chain homotopy after transfer to
$K\vee2S^5$, which has no six-cells. Tensoring with the unchanged
parity factor preserves this zero homotopy.

The additional tensor-zero contraction has the same property
for a functorial reason. It is obtained by first applying the
coefficient retraction $Y\to K$, then the contraction in
$T_K(M)$, and then the coefficient section. On relative
coefficient chains the first map is the zero map from the two
five-sphere cells. The relative tensor chain functor is the
tensor product of this coefficient complex and the reduced
parity complex; hence the entire resulting chain homotopy is
zero on that model. The contraction's naturality, rather than
merely the acyclicity of its intermediate target, justifies
this assertion. Base collars factor through the base and
likewise contribute nothing to the relative complex.

Renumber degree seven as zero and put $E=R^2$. The first cone
has $K_n=E\oplus E$ for $n>0$, $K_0=E$, and
$d(a,b)=(f_{n-1}b,0)$. The lower coordinate of $\ell_n$ is
$1-f_{n-1}=f_n$, and its upper coordinate is zero.
The second summand has upper coordinate
$\epsilon_n(1-f_n)$ and lower coordinate zero, where
$\epsilon_n=\pm1$ includes the parity-shift and pinch
orientations. Therefore $g$ induces the isomorphism
\[
 E\longrightarrow H_n(K)=(1-f_n)E\oplus f_nE,\qquad
 x\longmapsto(\epsilon_n(1-f_n)x,f_nx)\quad(n>0),
\]
with inverse $(u,v)\mapsto\epsilon_nu+v$.
Here is an explicit chain deformation retraction, so this is
stronger than a homology calculation. Write $p_n=f_n$ and
$q_n=1-f_n$, so $p_{n-1}=q_n$. Project the first cone onto its
zero-differential homology summand by
$\pi_n(a,b)=(q_na,p_nb)$, with $H_0=q_0E$.
The evident inclusion $\iota$ and the degree-one map
\[
 s_n(a,b)=(0,p_na)\in K_{n+1}
\]
satisfy $ds+sd=1-\iota\pi$, including the bottom degree.
The attaching map lands in this chosen homology summand.
Its cone therefore splits into contractible two-term pairs in
every positive degree and leaves exactly $(1-e)E$ at the bottom.
Restoring the shift gives the claimed contraction onto
$(1-e)E[7]$. Independent signs on both coordinates would have
the same proof, using inverse $(u,v)\mapsto\alpha_nu+\beta_nv$.

All relative pieces are simply connected above the base, and
the two cones apply van Kampen to maps inducing the identity
on $\pi_1K$. Thus $\pi_1(B\mathcal W_8)=Q_8$.
The finite base has universal-cover free chain ranks
$(10,58,96,48)$ in degrees zero through three. The connecting
chain map from the relative projective module in degree seven
to its suspension vanishes for dimension reasons. Consequently
\[
 C_*(\widetilde{B\mathcal W_8})\simeq
 C_*(\widetilde{BK_8})\oplus(1-e)R^2[7]
\]
over $R$. For completeness there is a direct domination argument.
The nerve has finitely many cells in each degree, so its
seven-skeleton $L_7$ is a finite CW complex. The inclusion into
$B\mathcal W_8$ is seven-connected. Replace it by a fibration;
its fibre is six-connected. Every obstruction to a homotopy
section lies in $H^j(B\mathcal W_8;\mathcal A)$ for a local
coefficient system $\mathcal A$ and $j\geq8$.
These groups all vanish by the bounded projective chain model
just computed. Inductive obstruction theory gives a homotopy
section, so $L_7$ dominates $B\mathcal W_8$.
The mapping telescope of the resulting homotopy idempotent
on $L_7$ is a CW model of dimension at most eight.
This also verifies directly the hypotheses in Wall's
finite-domination criterion \cite{wall1965D}.
Its reduced projective Euler class is
\[
 (-1)^7[(1-e)R^2]=[eR^2]=[\operatorname{im}e_8]\ne0.
\]
Wall's obstruction theorem \cite[Theorem F]{wall1965D}
excludes a finite CW model.

\paragraph{Verification and exact size.}
The ordinal, product, span, and height-one composition rules
prove all category laws universally. Complete small tables
independently check those formulas. The nary flattening audit
checks 6\,765\,548 simplex/pullback equations, 3424 endpoint
equations, and 12\,280 inverse equations in a nontrivial small
instance. The actual $Q_8$ coefficient paths check all 64 group
disks and 512 group/sign/index cases. The final run checks 558
base-retraction equations, 468 base-transformation equations,
and 15\,588 identity/associativity equations at the two integer
ends of every final type block. The last checks sample the
enormous table; they are not an exhaustive enumeration.
The exact type-count recursion gives the counts in the theorem,
and also
\[
\begin{split}
 \#\{\text{typed pairs}\}&=
 138\,512\,978\,818\,954\,605\,575\,263\,121\,838,\\
 \#\{\text{typed triples}\}&=
 52\,322\,049\,455\,781\,789\,059\,845\,756\,837\,400\,964.
\end{split}
\]
The machine-readable record is
\texttt{outputs/d\_final\_wall\_category\_summary.json}.
\end{proof}
